\subsection{可分裂半单李代数的自同构}

命题 5. 假设 g 是可分裂的。群 \(\operatorname{Aut}_{0}(\mathfrak{g})\) 单可递地作用在 g 的标架组成的集合上。

设 \(e_1 = (\mathfrak{g}, \mathfrak{h}_1, \mathrm{B}_1, (X_\alpha^1)_{\alpha \in \mathrm{B}_1})\) ，\(e_2 = (\mathfrak{g}, \mathfrak{h}_2, \mathrm{B}_2, (X_\alpha^2)_{\alpha \in \mathrm{B}_2})\) 是 \(\mathfrak{g}\) 的两个标架。存在 \(\operatorname{Aut}_0(\mathfrak{g})\) 的至少一个元素把 \(e_1\) 变为 \(e_2\) （命题 1 与命题 4）。设 \(\bar{k}\) 是 \(k\) 的一个代数闭包。存在 \(\operatorname{Aut}_e(\mathfrak{g} \otimes_k \bar{k})\) 的一个元素把 \(\mathfrak{h}_1 \otimes_k \bar{k}\) 变为 \(\mathfrak{h}_2 \otimes_k \bar{k}\) （第七章 §3 第 2 节 定理 1）。于是，由命题 4 与命题 2 的推论 2，存在元素 \(\varphi\) （\(\operatorname{Aut}_e(\mathfrak{g} \otimes_k \bar{k})\) 的），它把标架 \((\mathfrak{g} \otimes_k \bar{k}, \mathfrak{h}_1 \otimes_k \bar{k}, \mathrm{B}_1, (X_\alpha^1)_{\alpha \in \mathrm{B}_1})\) （\(\mathfrak{g} \otimes_k \bar{k}\) 的）变为标架 \((\mathfrak{g} \otimes_k \bar{k}, \mathfrak{h}_2 \otimes_k \bar{k}, \mathrm{B}_2, (X_\alpha^2)_{\alpha \in \mathrm{B}_2})\) 。由于 \(\mathfrak{h}_1\) 与诸 \(X_\alpha^1\) （相应地 \(\mathfrak{h}_2\) 与诸 \(X_\alpha^2\) ）生成 \(\mathfrak{g}_1\) （相应地 \(\mathfrak{g}_2\) ），我们有 \(\varphi(\mathfrak{g}_1) = \mathfrak{g}_2\) ，故 \(\varphi\) 形如 \(\psi \otimes 1\) ，其中 \(\psi \in \operatorname{Aut}_0(\mathfrak{g})\) ，且 \(\psi\) 把 \(e_1\) 变为 \(e_2\) 。

推论 1. 设 \((\mathfrak{g},\mathfrak{h},\mathrm{B},(X_{\alpha})_{\alpha\in\mathrm{B}})\) 是 g 的一个标架，G 是使该标架不变的 g 的自同构组成的群（同构于 \(\operatorname{Aut}(\mathrm{R},\mathrm{B})\) ）。则 \(\operatorname{Aut}(\mathfrak{g})\) 是 G 与 \(\operatorname{Aut}_{0}(\mathfrak{g})\) 的半直积。

事实上，\(\operatorname{Aut}(\mathfrak{g})\) 的每个元素把 \((\mathfrak{g},\mathfrak{h},\mathrm{B},(X_{\alpha})_{\alpha\in\mathrm{B}})\) 变为 g 的一个标架。由命题 5，\(\operatorname{Aut}(\mathfrak{g})\) 模 \(\operatorname{Aut}_{0}(\mathfrak{g})\) 的每个陪集与 G 恰好交于一点。证毕。

由推论 1 可知，群 \(\operatorname{Aut}(\mathfrak{g})/\operatorname{Aut}_{0}(\mathfrak{g})\) 可以与 \(\operatorname{Aut}(\mathrm{R},\mathrm{B})\) 等同，并且同构于 R 的邓金图的自同构群。

推论 2. \(\operatorname{Aut}(\mathfrak{g}) = \operatorname{Aut}_0(\mathfrak{g})\) ，当 \(\mathfrak{g}\) 是 \(\mathrm{A}_1\) , \(\mathrm{B}_n\) ( \(n \geq 2\) ), \(\mathrm{C}_n\) ( \(n \geq 2\) ), \(\mathrm{E}_7\) , \(\mathrm{E}_8\) , \(\mathrm{F}_4\) , \(\mathrm{G}_2\) 型的可分裂单李代数时。商群 \(\operatorname{Aut}(\mathfrak{g}) / \operatorname{Aut}_0(\mathfrak{g})\) 的阶是 2 ，当 \(\mathfrak{g}\) 是 \(\mathrm{A}_n\) ( \(n \geq 2\) ), \(\mathrm{D}_n\) ( \(n \geq 5\) ), \(\mathrm{E}_6\) 型时；它同构于 \(\mathfrak{S}_3\) ，当 \(\mathfrak{g}\) 是 \(\mathrm{D}_4\) 型时。

这由推论 1 与第六章图版 I 至 IX 得出。

注. 1) 设 \(e_1 = (\mathfrak{g}, \mathfrak{h}_1, \mathrm{B}_1, (X_\alpha^1)_{\alpha \in \mathrm{B}_1})\) ，\(e_2 = (\mathfrak{g}, \mathfrak{h}_2, \mathrm{B}_2, (X_\alpha^2)_{\alpha \in \mathrm{B}_2})\) ，\(e_2' = (\mathfrak{g}, \mathfrak{h}_2, \mathrm{B}_2, (Y_\alpha^2)_{\alpha \in \mathrm{B}_2})\) 是 \(\mathfrak{g}\) 的标架，而 \(s\) （相应地 \(s'\) ）是 \(\operatorname{Aut}_0(\mathfrak{g})\) 中把 \(e_1\) 变为 \(e_2\) （相应地 \(e_2'\) ）的元素。则 \(s|_{\mathfrak{h}_1} = s'|_{\mathfrak{h}_1}\) 。事实上，\(s'^{-1}s \in \operatorname{Aut}_0(\mathfrak{g}, \mathfrak{h}_1)\) 且 \(s'^{-1}s(\mathrm{B}_1) = \mathrm{B}_1\) ，故 \(\varepsilon(s'^{-1}s) = 1\) 。

\begin{enumerate}
\def\labelenumi{\arabic{enumi})}
\setcounter{enumi}{1}
\tightlist
\item
  设 X 是偶 \((\mathfrak{h},\mathrm{B})\) 组成的集合，其中 h 是 g 的分裂嘉当子代数，B 是 \((\mathfrak{g},\mathfrak{h})\) 的根系的一个基。若 \(x=(\mathfrak{h},\mathrm{B})\) 与 \(x'=(\mathfrak{h}',\mathrm{B}')\) 是 X 的两个元素，则存在 \(s\in\operatorname{Aut}_{0}(\mathfrak{g})\) 把 x 变为 \(x'\) （命题 5），且 s 在 h 上的限制 \(s_{x',x}\) 与 s 的选取无关（注 1）。特别地，\(s_{x'',x'}\circ s_{x',x}=s_{x'',x}\) 若 \(x,x',x''\in X\) ，且 \(s_{x,x}=1\) 。满足条件的诸族 \((h_{x})_{x\in\mathrm{X}}\) 组成的集合
\end{enumerate}

\[
a) h _ {x} \in \mathfrak {h} \text {   if   } x = (\mathfrak {h}, B)
\]

\[
b) s _ {x ^ {\prime}, x} (h _ {x}) = h _ {x ^ {\prime}} \text { if } x, x ^ {\prime} \in X
\]

该集合以自然的方式成为一个向量空间 \(\mathfrak{h}(\mathfrak{g})\) ，我们有时称它为 g 的典则嘉当子代数。对 \(x = (\mathfrak{h}, \mathrm{B})\) 与 \(x' = (\mathfrak{h}', \mathrm{B}')\) ，\(s_{x', x}\) 把 B 变为 \(B'\) ，从而把 \((\mathfrak{g}, \mathfrak{h})\) 的根系变为 \((\mathfrak{g}, \mathfrak{h}')\) 的根系；由此，\(\mathfrak{h}(\mathfrak{g})^*\) （\(\mathfrak{h}(\mathfrak{g})\) 的对偶）自然地配备了一个根系 \(\mathrm{R}(\mathfrak{g})\) 以及一个基 \(\mathrm{B}(\mathfrak{g})\) （\(\mathrm{R}(\mathfrak{g})\) 的基）。我们有时称 \(\mathrm{R}(\mathfrak{g})\) 为 g 的典则根系，称 \(\mathrm{B}(\mathfrak{g})\) 为它的典则基。群 \(\mathrm{Aut}(\mathfrak{g})\) 作用在 \(\mathfrak{h}(\mathfrak{g})\) 上，使 \(\mathrm{R}(\mathfrak{g})\) 与 \(\mathrm{B}(\mathfrak{g})\) 保持稳定；\(\mathrm{Aut}(\mathfrak{g})\) 中平凡地作用在 \(\mathfrak{h}(\mathfrak{g})\) 上的元素正是 \(\mathrm{Aut}_0(\mathfrak{g})\) 的元素。

命题 6. 设 \(\mathfrak{h}\) 是 \(\mathfrak{g}\) 的一个分裂嘉当子代数。以第 1 节的记号，我们有 \(\mathrm{Aut}_0(\mathfrak{g}) = \mathrm{Aut}_e(\mathfrak{g}).\mathrm{Ker}\varepsilon = \mathrm{Aut}_e(\mathfrak{g}).f(\mathrm{T}_Q)\) 。

由 §3 第 3 节 命题 10 的推论，\(\mathrm{Aut}_0(\mathfrak{g}) = \mathrm{Aut}_e(\mathfrak{g}).\mathrm{Aut}_0(\mathfrak{g},\mathfrak{h})\) 。另一方面，由 §2 第 2 节 定理 2 的推论，\(\varepsilon (\mathrm{Aut}_e(\mathfrak{g},\mathfrak{h}))\supset \mathrm{W}(\mathrm{R})\) ，故 \(\mathrm{Aut}_0(\mathfrak{g},\mathfrak{h}) = \mathrm{Aut}_e(\mathfrak{g},\mathfrak{h}).\mathrm{Ker}\varepsilon\) 。

注 3. 命题 6 表明，典则同态

\[
\iota : \mathrm{T} _ {\mathrm{Q}} / \operatorname{Im} (\mathrm{T} _ {\mathrm{P}}) \to \operatorname{Aut} _ {0} (\mathfrak {g}) / \operatorname{Aut} _ {e} (\mathfrak {g}),
\]

（由第 2 节的图诱导的）是满射。特别地，\(\mathrm{Aut}_e(\mathfrak{g})\) 包含 \(\mathrm{Aut}_0(\mathfrak{g})\) 的换位子群；我们将看到（§11 第 2 节 命题 3）二者实际上相等。此外，可以证明 \(\iota\) 是单射，换言之

\[
f (\mathrm{T} _ {\mathrm{Q}}) \cap \operatorname{Aut} _ {e} (\mathfrak {g}) = f ^ {\prime} (\mathrm{T} _ {\mathrm{P}}),
\]

（参见 §7 习题 26 d)）。

命题 7. 设 g 是一个可分裂半单李代数，b 是 g 的一个 Borel 子代数，\(p_{1}\) 与 \(p_{2}\) 是包含 b 的两个不同的抛物子代数。则 \(p_{1}\) 与 \(p_{2}\) 在 \(\operatorname{Aut}_{0}(\mathfrak{g})\) 下不共轭。

我们可以假设 \(k\) 是代数闭的。设 \(s \in \mathrm{Aut}_0(\mathfrak{g})\) 满足 \(s(\mathfrak{p}_1) = \mathfrak{p}_2\) 。设 \(\mathfrak{h}\) 是 \(\mathfrak{g}\) 的一个包含在 \(\mathfrak{b} \cap s(\mathfrak{b})\) 中的嘉当子代数（§3 第 3 节 命题 10）。由于 \(\mathfrak{h}\) 与 \(s(\mathfrak{h})\) 都是 \(s(\mathfrak{b})\) 的嘉当子代数，存在 \(u \in [\mathfrak{b}, \mathfrak{b}]\) 使 \(e^{\mathrm{ad}u}(\mathfrak{h}) = s(\mathfrak{h})\) （第七章 §3 第 4 节 定理 3）。把 \(s\) 换成 \(e^{-\mathrm{ad}u}s\) ，我们化归到 \(s(\mathfrak{h}) = \mathfrak{h}\) 的情形，于是 \(s\) 在 \(\mathfrak{h}\) 上诱导出 Weyl 群 W 的一个元素 \(\sigma\) （W 是 \((\mathfrak{g}, \mathfrak{h})\) 的 Weyl 群）（命题 4）。设 C 是对应于 \(\mathfrak{b}\) 的 Weyl 房。则 \(\mathfrak{p}_1\) 与 \(\mathfrak{p}_2\) 对应于面 \(F_1\) 与 \(F_2\) （\(\mathfrak{h}_{\mathbf{R}}\) 的包含在 C 的闭包中的面）。我们有 \(\sigma(F_1) = F_2\) 。由于 \(\sigma \in W\) ，这蕴含 \(F_1 = F_2\) （第五章 §3 第 3 节 定理 2），故 \(\mathfrak{p}_1 = \mathfrak{p}_2\) 。

注 4. 设 g 是一个可分裂半单李代数，P 是 g 的抛物子代数组成的集合，\(\operatorname{Aut}_{0}(\mathfrak{g})\) 作用在这个集合上。沿用注 2 的记号。设 \(\Sigma\) 是 \(\mathrm{B}(\mathfrak{g})\) 的一个子集。给定 \(\Sigma\) 等价于对每个 \(x = (\mathfrak{h}, \mathrm{B}) \in \mathrm{X}\) 给定 B 的一个子集 \(\Sigma_{x}\) ，使 \(s_{x',x}\) 把 \(\Sigma_{x}\) 变为 \(\Sigma_{x'}\) 对任何 \(x, x' \in X\) 成立。设 \(p_{x}\) 是对应于 \(\Sigma_{x}\) 的 g 的抛物子代数（§3 第 4 节 注）。\(p_{x}\) 在 \(\operatorname{Aut}_{0}(\mathfrak{g})\) 作用下的轨道是诸 \(p_{x'}\) （\(x' \in X\) ）组成的集合。这就定义了一个从 \(\mathfrak{P}(\mathrm{B}(\mathfrak{g}))\) 到 \(\mathcal{P}/\operatorname{Aut}_{0}(\mathfrak{g})\) 的映射。这个映射由 §3 第 4 节的注知是满射，由命题 7 知是单射。

